\chapter{Hipótesis y Objetivos}
\label{ch:obj}

% TODO (tesis): hipótesis y OE se conservan tal cual (son el contrato del
% anteproyecto); la sección "Tareas asociadas" puede eliminarse o moverse a
% apéndice. OJO: los OE y alcances mencionan las 6 categorías del esquema
% viejo (Arithmetic/Logic/...); decidir con el asesor cómo documentar el
% cambio al esquema por firma de energía sin romper el contrato de los OE.

% -----------------------------------------------------------------------
\section{Hipótesis}
\label{sec:hipotesis}

Un módulo hardware de clasificación de instrucciones integrado en el procesador
CV32E40P, que expone conteos por categoría mediante registros CSR, permite
estimar el consumo energético de aplicaciones bare-metal con un error relativo
medio inferior al 10\,\% respecto a la medición eléctrica directa sobre la
plataforma FPGA Nexys~A7-100T.

El umbral del 10\,\% se adopta como criterio de desempeño para evaluar la
utilidad del modelo en la comparación de regiones de código. Este margen toma
en cuenta que la plataforma objetivo opera en bare-metal sobre FPGA, sin
infraestructura de monitoreo energético del sistema operativo, y que cada
categoría agrupa instrucciones con comportamientos eléctricos distintos. Bajo
estas condiciones, el objetivo no es reemplazar la medición eléctrica de
referencia durante la caracterización, sino obtener una estimación suficientemente
precisa para orientar el análisis energético del firmware.

% -----------------------------------------------------------------------
\section{Objetivo general}
\label{sec:obj_general}

Diseñar e implementar un módulo hardware de clasificación de instrucciones
integrado en el procesador CV32E40P sobre la plataforma PULPissimo, que genere
conteos por categoría para estimar, mediante un flujo externo de análisis, el
consumo energético de aplicaciones bare-metal ejecutadas en la FPGA.

% -----------------------------------------------------------------------
\section{Objetivos específicos}
\label{sec:obj_especificos}

\begin{enumerate}

    \item Diseñar y construir el circuito de medición de potencia basado en
          resistor de sensado de corriente y amplificador de sensado INA240
          en configuración \emph{high-side}, y verificar su correcto
          funcionamiento sobre la línea de alimentación de la placa Nexys~A7-100T.

    \item Implementar en SystemVerilog el módulo clasificador de instrucciones
          como bloque RTL integrado en el núcleo CV32E40P, exponiendo un
          contador de 64 bits por cada una de las seis categorías de
          instrucción definidas mediante pares de CSR de 32 bits.

    \item Sintetizar el SoC PULPissimo con el módulo clasificador incorporado
          sobre la FPGA Nexys~A7-100T y verificar el funcionamiento de los
          registros CSR mediante firmware bare-metal y la cadena JTAG.

    \item Caracterizar experimentalmente los coeficientes de energía $e_i$
          por categoría de instrucción sobre la plataforma FPGA, aplicando
          el método de bucles dominados por categoría y el método de regresión
          por histograma.

    \item Desarrollar el pipeline de análisis de datos que automatice el
          procesamiento de las mediciones de potencia adquiridas por el sistema
          ADS1115--LilyGO ESP32 y la lectura de registros CSR, y genere
          estimaciones de energía, potencia promedio y desglose por categoría.

    \item Validar el modelo de estimación comparando la potencia promedio estimada
          $\bar{P}_{\text{est}} = E_{\text{est}}/\Delta t$ contra la potencia promedio medida
          físicamente en la plataforma FPGA durante la ejecución de cargas de
          trabajo mixtas no utilizadas en la caracterización.

\end{enumerate}

% -----------------------------------------------------------------------
\section{Tareas asociadas a los objetivos específicos}
\label{sec:tareas_objetivos}

\begin{enumerate}[label=\textbf{OE\arabic*:}, leftmargin=*]

    \item Circuito de medición de potencia.
    \begin{compactitem}
        \item Establecer la topología de medición \emph{high-side} sobre la línea
              de alimentación de la Nexys~A7-100T.
        \item Dimensionar el resistor de sensado a partir del rango de corriente
              definido para la plataforma.
        \item Conectar y probar la cadena INA240A1--ADS1115--LilyGO ESP32.
        \item Calibrar la conversión de tensión medida a corriente y potencia.
    \end{compactitem}

    \item Módulo clasificador de instrucciones.
    \begin{compactitem}
        \item Establecer y justificar las categorías de instrucciones a
              contabilizar.
        \item Seleccionar y justificar las señales del pipeline del CV32E40P
              utilizadas para detectar instrucciones retiradas válidamente.
        \item Implementar la lógica de clasificación en SystemVerilog.
        \item Implementar los contadores de 64 bits por categoría y su acceso
              mediante pares de CSR de 32 bits.
    \end{compactitem}

    \item Integración y verificación en FPGA.
    \begin{compactitem}
        \item Integrar el clasificador con el banco de CSR del CV32E40P.
        \item Sintetizar el SoC PULPissimo modificado para la FPGA Nexys~A7-100T.
        \item Cargar y ejecutar firmware bare-metal mediante la cadena
              GDB--OpenOCD--JTAG.
        \item Verificar la lectura y el reinicio de los contadores mediante acceso CSR.
    \end{compactitem}

    \item Caracterización de coeficientes energéticos.
    \begin{compactitem}
        \item Implementar rutinas dominadas por categoría para las categorías
              de instrucción.
        \item Medir la potencia promedio de cada categoría en régimen estable.
        \item Ejecutar cargas mixtas y recolectar histogramas de instrucciones.
        \item Calcular coeficientes mediante bucles dominados por categoría y
              regresión por histograma.
    \end{compactitem}

    \item Pipeline de análisis de datos.
    \begin{compactitem}
        \item Automatizar la lectura de mediciones adquiridas por el sistema
              ADS1115--LilyGO ESP32.
        \item Procesar las muestras de potencia y calcular energía por ejecución.
        \item Integrar la lectura de contadores CSR en el flujo de análisis.
        \item Generar estimaciones, tablas y visualizaciones por categoría.
    \end{compactitem}

    \item Validación del modelo.
    \begin{compactitem}
        \item Definir un conjunto de cargas mixtas no utilizadas durante la
              caracterización.
        \item Ejecutar las cargas sobre la plataforma FPGA y medir su potencia.
        \item Comparar la potencia promedio estimada contra la potencia promedio medida.
        \item Calcular el error relativo y verificar el cumplimiento de la
              hipótesis.
    \end{compactitem}

\end{enumerate}

% -----------------------------------------------------------------------
\section{Alcances y limitaciones}
\label{sec:alcances_limitaciones}

El proyecto abarca el diseño, implementación y validación experimental de un
módulo hardware de clasificación de instrucciones integrado en el procesador
CV32E40P dentro del SoC PULPissimo. La implementación se realizará sobre la
plataforma FPGA Nexys~A7-100T y estará orientada a firmware bare-metal, sin
dependencia de sistema operativo.

La estimación energética se limitará a un modelo por categoría de instrucción.
Para ello se considerarán seis categorías principales: \emph{Arithmetic},
\emph{Logic}, \emph{Memory}, \emph{Branch}, \emph{Jump} y
\emph{Floating point}. Cada categoría tendrá un contador de 64 bits expuesto mediante pares de CSR
personalizados de 32 bits.

El trabajo incluye la construcción y uso de una cadena de medición de potencia
externa basada en un sensor de corriente INA240A1, un convertidor ADS1115 y una
tarjeta ESP32 para adquisición de datos. Esta cadena se utilizará como
referencia física para caracterizar los coeficientes energéticos del modelo y
validar la estimación calculada a partir de los contadores internos.

La caracterización experimental considerará tanto el método de bucles
dominados por categoría como el método de regresión por histograma. Los
coeficientes obtenidos por ambos métodos se compararán según su error frente a mediciones
físicas de referencia, y el conjunto con mejor desempeño se utilizará para la
validación final del modelo.

Como limitación principal, los coeficientes energéticos obtenidos serán válidos
para la configuración experimental utilizada: CV32E40P integrado en PULPissimo,
sintetizado sobre la FPGA Nexys~A7-100T y operando bajo las condiciones de
alimentación, frecuencia y bitstream definidas durante la caracterización. Por
tanto, no se asume que los coeficientes puedan trasladarse directamente a un
ASIC, a otra FPGA o a otra frecuencia de operación.

La medición eléctrica corresponde al consumo observado en la plataforma FPGA
completa bajo la configuración experimental, no al consumo aislado del núcleo
CV32E40P como circuito independiente. El modelo resultante estima el consumo
asociado a una ejecución en esa plataforma, pero no separa
explícitamente el aporte de todos los bloques internos del SoC ni de los
reguladores de la placa.

El clasificador trabaja a nivel de categorías de instrucción, no a nivel de
instrucciones individuales. En consecuencia, instrucciones distintas dentro de
una misma categoría comparten un coeficiente energético promedio, por lo que el
modelo puede perder precisión en programas donde una categoría agrupe
operaciones con comportamientos eléctricos muy diferentes.

La validación se realizará con cargas bare-metal controladas y no contempla
ejecución bajo sistema operativo, multiprogramación, interrupciones complejas ni
periféricos externos con patrones de actividad no controlados. Además, el modelo
no pretende reemplazar la medición eléctrica durante la caracterización inicial;
la instrumentación externa sigue siendo necesaria para obtener los coeficientes
y validar la hipótesis.

%%% Local Variables:
%%% mode: latex
%%% TeX-master: "main"
%%% End:
